The lower ends of the brackets come from a county-aware recombination hill climber followed by exact large-neighbourhood search; to see what certification adds we compare them with the incumbent of a plain unconstrained ReCom short-burst run of 40~s per (state, party, $q$) with the same population tolerance (circles in Figure~\ref{fig:ablation}). This is a modest heuristic budget compared with the hours of computation behind published maps \citep{atlas2026}, so the comparison is illustrative rather than a ranking of methods.

\paragraph{Heuristics undershoot.}
The verified plans found under the budget (which is a restriction!) have a higher $q$-th margin than the unconstrained incumbent in 45 of 70 comparisons (median $+0.09$ point; range $-0.9$ to $+2.1$). In the largest cases the incumbent is one to two points below what is provably attainable: West Virginia Republicans 74.0\% against a certified 76.0--77.0\%, Idaho Republicans 74.5\% against 75.5--76.0\%.

\paragraph{A certified price of county integrity.}
In one case the unconstrained incumbent \emph{exceeds} the certified upper end: Maine Democrats reach 63.3\% in a verified unconstrained plan, whereas no plan of the budgeted regime reaches 63.0\%. Together the two facts are a certificate that the budget of one county split reduces the best-district capacity of Maine Democrats by at least 0.3 point at this tolerance. For every other state and party in our data the incumbent lies inside or below the bracket, so no separation is possible without a stronger unconstrained search or a certified upper bound for the unconstrained regime, which the negative result above suggests is out of reach at this resolution.
